% !TeX root = ../main.tex
\section{讨论}
\subsection{三种思路的连续性与差异}
V2、V3、V4 分别回答了不同层次的问题。V2 建立实际计算及其证据的可复核性；V3 研究在有限任务范围内，如何将自然语言语义、工具组合和失败修复连接起来；V4 则进一步考察任务算法本身能否由代码生成，并在受控边界内完成开放目标。三者之间存在问题上的递进，但并非在共同任务和相同条件下进行的性能迭代。

V3 的 18/18 包含固定任务、脚本化对话、自动确认与拒绝情境；V4 的计算成功指标针对实际 ORCA 调用，冻结批次为 7/8，后续案例还覆盖更多任务类型。两类分母不同，不能直接解释为性能下降或提升。若要比较版本收益，需要相同原始目标、输入、最低证据、模型与资源，并分别测量完整完成、数值误差和成本。

\subsection{V4 的主要瓶颈是证据链的一致性}
V4 已通过生成程序获得优化结构、振动频率、有限构象单点能量、激发能及核磁屏蔽等真实结果，证明该路线可以支持多种计算任务。当前困难主要出现在从成功计算到后续行动的连接处。失败可归纳为四类：语义与结构化表示不一致、目标条件在层间传递不完整、修复后目标与产物关联丢失，以及反馈上下文和错误发现不足。它们不属于传统 SCF 收敛问题，却可以决定整个科学任务是否完成。

这一结果将约束机制本身纳入研究对象。约束过弱可能使错误目标获得形式通过，过度重复或互不一致的约束又会阻断合法任务。有效改进应在保留科学要求、原始目标与资源边界的同时，减少同一事实的重复表示，并为当前失败提供直接可查询的证据。相关机制的修改需要新的冻结评测；本文新增的计算成功标记仅明确原始 ORCA 执行事实，不改写历史工作流状态。

\subsection{闭环性与科学可靠性}
AutoDFT 为成功与失败均进入后续规划提供了方法参照，多参考自主计算研究则强调电子态身份和覆盖率的重要性。V4 将这些问题组织在单一决策循环中，以统一工具边界约束实际操作。这是架构选择，本文没有单决策者与多角色系统的受控比较，不能据此宣称性能优势。

闭环性也不能仅由存在反馈接口推断。真正有效的反馈应改变后续行动，并持续缩小原始目标的证据缺口。T3 和 H1 均已成功执行 ORCA，同时表明后续反馈组织仍有完善空间。未来评价应考察结果是否被正确消费、必要操作是否实际执行，以及停止时的部分结论是否与已有证据一致。

\subsection{有效性威胁与进一步验证}
实验规模和异质性限制了统计推断。V2 为人工批准的选定成功链，V3 的任务与对话高度限定，V4 主评测每题仅一次；跨版本和跨配置数据不能组成统一成功概率。早期 V3 实验未保存完整执行源码快照，限制严格复现；部分 V4 案例具有人工协助，应与自主评测区分。

科学有效性同样具有范围限制。局部频率检查不证明全局构象最优，候选子集能量排名不证明自由能排序，原始张量存在也不证明目标已完成。固定设置下的数值重复性支持可重复计算，不替代独立高精度或实验参照。本文未验证 V4 的多参考求解能力，也未将指定隔离配置的测试扩展为普遍安全保证。

进一步研究应首先建立同配置的重复试验和未暴露任务集，再对科学输出引入独立参考或确定性复核。针对 V4，可分别检验目标条件统一、修复关联保持和反馈压缩的贡献；扫描任务需要检查端点与实际几何，跨环境比较需要检查条件和共同几何，张量任务需要检查单位及坐标变换。由此将“可以生成程序”的证据推进为“可以稳定完成科学目标”的证据。
